% ============================================================
%  HalluGuard: Real-Time Confidence-Adaptive Hallucination
%  Detection and Mitigation in Retrieval-Augmented Generation
%  without Fine-Tuning
% ============================================================
\documentclass[conference]{IEEEtran}
\IEEEoverridecommandlockouts

\usepackage{cite}
\usepackage{amsmath,amssymb,amsfonts}
\usepackage{algorithmic}
\usepackage{algorithm}
\usepackage{graphicx}
\usepackage{textcomp}
\usepackage{xcolor}
\usepackage{booktabs}
\usepackage{multirow}
\usepackage{url}
\usepackage{hyperref}
\usepackage{tikz}
\usetikzlibrary{shapes.geometric, arrows.meta, positioning, fit}

\graphicspath{{figures/}}

\def\BibTeX{{\rm B\kern-.05em{\sc i\kern-.025em b}\kern-.08em
    T\kern-.1667em\lower.7ex\hbox{E}\kern-.125emX}}

\begin{document}

\title{HalluGuard: Real-Time Confidence-Adaptive Hallucination Detection and Mitigation in Retrieval-Augmented Generation without Fine-Tuning}

\author{\IEEEauthorblockN{Arnab Mondal}
\IEEEauthorblockA{\textit{Dept.\ of Computer Science and Engineering (AI/ML)} \\
\textit{CHRIST (Deemed to be University)}\\
Bengaluru, India\\
arnab.mondal@mtech.christuniversity.in}
\and
\IEEEauthorblockN{Julian Benadit P}
\IEEEauthorblockA{\textit{Dept.\ of Computer Science and Engineering (AI/ML)} \\
\textit{CHRIST (Deemed to be University)}\\
Bengaluru, India\\
julian.p@christuniversity.in}
}

\maketitle

% ============================================================
%  ABSTRACT
% ============================================================
\begin{abstract}
Retrieval-Augmented Generation (RAG) systems hallucinate on 15--25\% of queries despite retrieval grounding, yet existing approaches either detect hallucinations or mitigate them---not both---and typically require model fine-tuning. We present \textsc{HalluGuard}, a training-free framework that integrates an ensemble hallucination detector with a confidence-gated mitigation router. The detector fuses Natural Language Inference (NLI), embedding similarity, and LLM-judge signals via parallel execution. A novel confidence-adaptive router selects proportionate correction strategies---query rewriting, context reranking, or self-reflection---based on detection confidence, eliminating the one-size-fits-all limitation of prior corrective systems. On 500 HaluEval QA samples, HalluGuard reduces the hallucination rate from 58.2\% to 39.0\% and improves mean faithfulness from 0.4227 to 0.5712. Ablation experiments confirm that the adaptive router achieves +46.4\% faithfulness improvement over baseline, surpassing fixed query rewriting (+36.9\%) and fixed self-reflection (+39.2\%). The detector attains an AUROC of 0.748 without any fine-tuning and runs entirely on CPU hardware.
\end{abstract}

\begin{IEEEkeywords}
Retrieval-Augmented Generation, Hallucination Detection, Faithfulness Evaluation, Confidence-Adaptive Mitigation, Large Language Models
\end{IEEEkeywords}

% ============================================================
%  I. INTRODUCTION
% ============================================================
\section{Introduction}

Large Language Models (LLMs) are increasingly deployed in enterprise, healthcare, and legal domains where factual accuracy is paramount~\cite{ji2023survey}. Despite remarkable fluency, LLMs are prone to \textit{hallucination}---generating content that is unfaithful to provided evidence or factually incorrect~\cite{ji2023survey}. The consequences are severe: in a widely reported 2023 incident, a lawyer submitted a legal brief to a US federal court containing six fabricated case citations generated by ChatGPT, resulting in sanctions.

Retrieval-Augmented Generation (RAG) has emerged as the dominant paradigm for grounding LLM outputs in external knowledge~\cite{lewis2020rag}. By conditioning generation on retrieved passages, RAG systems reduce hallucination relative to closed-book models. However, RAG does not eliminate hallucination. Three distinct failure modes persist:
(1)~\textit{Faithfulness failure}, where the generator ignores retrieved context and produces unsupported claims;
(2)~\textit{Retrieval-induced hallucination}, where noisy or irrelevant passages mislead the generator; and
(3)~\textit{Extrinsic injection}, where the LLM supplements retrieved evidence with parametric knowledge that may be incorrect.

Existing detection methods such as Luna~\cite{ravi2024luna} and MiniCheck~\cite{tang2024minicheck} provide accurate hallucination scoring but offer no correction mechanism. Corrective systems such as Self-RAG~\cite{asai2024selfrag} and RAG-HAT~\cite{wu2024raghat} address mitigation but require fine-tuning and apply a single fixed correction strategy regardless of hallucination severity. No prior system simultaneously achieves real-time detection, training-free operation, detection-mitigation integration, and confidence-adaptive strategy selection.

We present \textsc{HalluGuard}, a framework that closes this gap. Our contributions are:
\begin{enumerate}
    \item A \textbf{training-free ensemble detector} combining NLI, embedding similarity, and LLM-judge signals with parallel execution via \texttt{ThreadPoolExecutor}.
    \item A \textbf{confidence-gated mitigation router}---the first system to select correction strategy (query rewriting, context reranking, or self-reflection) based on detection confidence.
    \item \textbf{Claim-level hallucination explanation} with per-sentence NLI scoring and natural language justification, enabling human-in-the-loop verification.
    \item \textbf{Multi-turn entity consistency tracking} that detects contradictions across conversation turns.
\end{enumerate}

% ============================================================
%  II. RELATED WORK
% ============================================================
\section{Related Work}

\subsection{NLI-Based Faithfulness Detection}

Lewis et al.~\cite{lewis2020rag} introduced Retrieval-Augmented Generation, establishing the paradigm of grounding LLM outputs in retrieved passages. Subsequent work on faithfulness evaluation leveraged Natural Language Inference (NLI) models to assess entailment between context and generated claims. Honovich et al.~\cite{honovich2022true} proposed TRUE, a unified evaluation framework benchmarking NLI-based factual consistency across summarisation and dialogue tasks. Tang et al.~\cite{tang2024minicheck} developed MiniCheck, an efficient fact-checking system that decomposes responses into atomic claims and verifies each against grounding documents, achieving strong performance with reduced model sizes. More recently, Kovacs et al.~\cite{kovacs2025lettucedetect} introduced LettuceDetect, a token-level hallucination detection framework specifically designed for RAG applications. While these methods achieve accurate detection, they are strictly evaluation tools that provide no corrective action upon detecting hallucination.

\subsection{Sampling-Based Consistency Methods}

An orthogonal line of work exploits the stochasticity of LLM generation to detect hallucination without reference documents. Manakul et al.~\cite{manakul2023selfcheckgpt} proposed SelfCheckGPT, which samples $k$ responses from the same model and identifies claims that appear inconsistently across samples as likely hallucinated. Farquhar et al.~\cite{farquhar2024semantic} extended this principle through Semantic Entropy, computing entropy over semantically clustered generations to detect confabulation. While reference-free, these methods require $k$ forward passes through the LLM (typically $k \geq 5$), making them prohibitively slow for real-time applications.

\subsection{Retrieval-Augmented Correction Systems}

Self-RAG~\cite{asai2024selfrag} trains an LLM to emit special reflection tokens that critique its own retrieval and generation decisions, enabling on-the-fly self-correction. However, Self-RAG requires supervised fine-tuning with reflection token annotations. Yan et al.~\cite{yan2024crag} proposed Corrective RAG (CRAG), which triggers web search when retrieval confidence is low, but applies a single fixed correction strategy. Wu et al.~\cite{wu2024raghat} introduced RAG-HAT, a hallucination-aware tuning pipeline that trains the generator to avoid hallucination, again requiring fine-tuning. A common limitation across these systems is the application of one fixed correction strategy regardless of hallucination severity.

\subsection{LLM-as-Judge Evaluation}

LLM-based evaluation has gained traction as a scalable alternative to human annotation. Es et al.~\cite{es2024ragas} introduced RAGAS, a framework that uses LLMs to evaluate RAG systems on faithfulness, answer relevance, and context recall. Min et al.~\cite{min2023factscore} proposed FActScore, which decomposes generated text into atomic facts and verifies each against a knowledge source. Sun et al.~\cite{sun2025redeep} developed ReDeEP, applying mechanistic interpretability to detect hallucination in RAG via internal model representations. While these methods exhibit high correlation with human judgement, they serve exclusively as evaluation tools with high per-query latency.

\medskip
\noindent\textbf{Positioning.} No prior system simultaneously achieves: (i)~real-time detection without multiple LLM forward passes, (ii)~training-free operation, (iii)~integrated detection and mitigation, and (iv)~confidence-adaptive strategy selection. HalluGuard is designed to fill precisely this gap.

% ============================================================
%  III. METHODOLOGY
% ============================================================
\section{Methodology}

Fig.~\ref{fig:architecture} presents the HalluGuard architecture. We describe each component below.

\begin{figure*}[t]
\centering
% TikZ Architecture Diagram for HalluGuard
% To be \input{} from halluguard_paper.tex
\tikzstyle{block} = [rectangle, draw, fill=blue!8, text centered,
    minimum height=2.4em, minimum width=6.5em, rounded corners=3pt,
    font=\small]
\tikzstyle{decision} = [diamond, draw, fill=orange!15, text centered,
    minimum height=2em, aspect=2.2, font=\small]
\tikzstyle{tier} = [rectangle, draw, fill=green!10, text centered,
    minimum height=2.2em, minimum width=7em, rounded corners=3pt,
    font=\small]
\tikzstyle{output} = [rectangle, draw, fill=gray!15, text centered,
    minimum height=2.2em, minimum width=6em, rounded corners=3pt,
    font=\small]
\tikzstyle{arrow} = [-{Stealth[length=2.5mm]}, thick]

\begin{tikzpicture}[node distance=1.1cm and 1.2cm]
    % Main pipeline (left to right)
    \node[block] (query) {User Query};
    \node[block, right=of query] (retriever) {Hybrid Retriever\\{\scriptsize Dense + BM25 + RRF}};
    \node[block, right=of retriever] (generator) {LLM Generator\\{\scriptsize Llama-3-8B}};
    \node[block, right=of generator] (detector) {Ensemble Detector\\{\scriptsize NLI + Emb + Judge}};
    \node[decision, right=1.4cm of detector] (decision) {$H > 0.5$?};

    % No branch (direct output)
    \node[output, above=1.0cm of decision] (output_no) {Final Response};

    % Yes branch — router tiers (below)
    \node[block, below=1.2cm of decision] (router) {Confidence Router};
    \node[tier, below left=0.9cm and 0.8cm of router] (tier1) {Query Rewrite\\{\scriptsize $0.5 < H < 0.65$}};
    \node[tier, below=0.9cm of router] (tier2) {Reranking\\{\scriptsize $0.65 \leq H < 0.80$}};
    \node[tier, below right=0.9cm and 0.8cm of router] (tier3) {Self-Reflection\\{\scriptsize $H \geq 0.80$}};

    % Corrected output
    \node[output, below=1.0cm of tier2] (output_yes) {Corrected Response};

    % Arrows — main flow
    \draw[arrow] (query) -- (retriever);
    \draw[arrow] (retriever) -- (generator);
    \draw[arrow] (generator) -- (detector);
    \draw[arrow] (detector) -- (decision);

    % Decision branches
    \draw[arrow] (decision) -- node[right, font=\scriptsize] {No} (output_no);
    \draw[arrow] (decision) -- node[right, font=\scriptsize] {Yes} (router);

    % Router to tiers
    \draw[arrow] (router) -- (tier1);
    \draw[arrow] (router) -- (tier2);
    \draw[arrow] (router) -- (tier3);

    % Tiers to output
    \draw[arrow] (tier1) -- (output_yes);
    \draw[arrow] (tier2) -- (output_yes);
    \draw[arrow] (tier3) -- (output_yes);

    % Explainer module (optional annotation)
    \node[block, fill=yellow!12, above=1.0cm of detector, minimum width=6em] (explainer) {Claim-Level\\Explainer};
    \draw[-{Stealth[length=2mm]}, dashed, thick, gray] (detector) -- (explainer);
\end{tikzpicture}

\caption{HalluGuard system architecture. A user query passes through hybrid retrieval and LLM generation, then enters the ensemble detector. Responses with hallucination score $H \leq 0.5$ are returned directly. Hallucinated responses are routed to one of three mitigation tiers based on confidence: query rewriting ($0.5 < H < 0.65$), context reranking ($0.65 \leq H < 0.80$), or self-reflection ($H \geq 0.80$).}
\label{fig:architecture}
\end{figure*}

\subsection{Hybrid Retriever}

The retrieval module combines dense and sparse retrieval to maximise recall. Dense retrieval uses sentence-transformers/all-MiniLM-L6-v2 embeddings stored in a ChromaDB vector store. Sparse retrieval uses BM25 via the rank-bm25 library. Results from both retrievers are fused using Reciprocal Rank Fusion (RRF)~\cite{lewis2020rag} with smoothing constant $k = 60$:
\begin{equation}
\text{RRF}(d) = \sum_{r \in \mathcal{R}} \frac{1}{k + \text{rank}_r(d)}
\label{eq:rrf}
\end{equation}
where $\mathcal{R}$ is the set of retrieval systems and $\text{rank}_r(d)$ is the rank of document $d$ in system $r$. The top-5 documents by RRF score form the context $\mathcal{C}$ passed to the generator.

\subsection{Ensemble Hallucination Detector}

The detector aggregates three complementary signals, each capturing a different dimension of faithfulness. All three signals execute in parallel via \texttt{ThreadPoolExecutor(max\_workers=3)}, reducing total detection latency from the sum to the maximum of individual signal latencies.

\medskip
\noindent\textbf{Signal 1: NLI Scorer.} We use the cross-encoder/nli-deberta-v3-base model to compute entailment probability for each sentence $s_i$ in the response against the concatenated context:
\begin{equation}
p_i = P(\text{entailment} \mid \text{premise}=\mathcal{C},\; \text{hypothesis}=s_i)
\end{equation}
\begin{equation}
S_{\text{NLI}} = 1 - \frac{1}{n}\sum_{i=1}^{n} p_i
\label{eq:nli}
\end{equation}
where $n$ is the number of sentences. Higher $S_{\text{NLI}}$ indicates lower faithfulness.

\medskip
\noindent\textbf{Signal 2: Embedding Similarity.} Using sentence-transformers/all-mpnet-base-v2, we compute cosine similarity between the response embedding and each context chunk embedding:
\begin{equation}
S_{\text{emb}} = 1 - \max_{c_j \in \mathcal{C}} \cos(\mathbf{e}_r, \mathbf{e}_{c_j})
\label{eq:emb}
\end{equation}
where $\mathbf{e}_r$ and $\mathbf{e}_{c_j}$ are the response and chunk embeddings respectively.

\medskip
\noindent\textbf{Signal 3: LLM Judge.} We prompt Llama-3-8B-Instruct (via Groq API) to produce a structured JSON assessment with a faithfulness score on a 1--5 scale, normalised to $[0, 1]$. The LLM judge provides complementary reasoning-based assessment and additionally returns problematic claims and natural language reasoning.

\medskip
\noindent\textbf{Weighted Fusion.} The final hallucination score combines all three signals:
\begin{equation}
H = w_1 \cdot S_{\text{NLI}} + w_2 \cdot S_{\text{emb}} + w_3 \cdot S_{\text{LLM}}
\label{eq:fusion}
\end{equation}
with weights $w_1 = 0.5$, $w_2 = 0.2$, $w_3 = 0.3$, reflecting the empirically observed reliability ordering NLI $>$ LLM judge $>$ embedding similarity. The detection decision is $\text{is\_hallucinated} = (H > 0.5)$.

\subsection{Confidence-Gated Mitigation Router}

The mitigation router is the primary contribution of this work. Unlike Self-RAG~\cite{asai2024selfrag} and CRAG~\cite{yan2024crag}, which apply a single fixed correction strategy, HalluGuard selects a strategy proportionate to the severity of detected hallucination. This design is motivated by the observation that mild hallucinations can be resolved with lightweight retrieval adjustments, while severe hallucinations require expensive iterative correction. Algorithm~\ref{alg:router} formalises the routing logic.

\begin{algorithm}[t]
\caption{Confidence-Gated Mitigation Router}
\label{alg:router}
\begin{algorithmic}[1]
\REQUIRE Query $q$, response $r$, context $\mathcal{C}$, score $H$
\ENSURE Corrected response $r^*$, strategy $\sigma$
\IF{$H \leq 0.5$}
    \STATE $\sigma \leftarrow$ \texttt{none}; \textbf{return} $r$, $\sigma$
\ELSIF{$0.5 < H < 0.65$}
    \STATE $\sigma \leftarrow$ \texttt{query\_rewrite}
    \STATE $Q' \leftarrow \textsc{RewriteQuery}(q, r)$ \COMMENT{3 alternative phrasings}
    \STATE $\mathcal{C}' \leftarrow \textsc{Retrieve}(Q')$ \COMMENT{re-retrieve with new queries}
    \STATE $\mathcal{C}' \leftarrow \textsc{Rerank}(q, \mathcal{C}', |\mathcal{C}|)$
    \STATE $r^* \leftarrow \textsc{Generate}(q, \mathcal{C}')$
\ELSIF{$0.65 \leq H < 0.80$}
    \STATE $\sigma \leftarrow$ \texttt{rerank}
    \STATE $\mathcal{C}' \leftarrow \textsc{Rerank}(q, \mathcal{C}, \text{top\_}k{=}3)$
    \STATE $r^* \leftarrow \textsc{Generate}(q, \mathcal{C}')$
\ELSE \COMMENT{$H \geq 0.80$}
    \STATE $\sigma \leftarrow$ \texttt{self\_reflection}
    \STATE $r^* \leftarrow \textsc{SelfReflect}(q, r, \mathcal{C}, \text{max\_iter}{=}2)$
\ENDIF
\STATE Score $r^*$; keep $r^*$ only if improved over $r$
\RETURN $r^*$, $\sigma$
\end{algorithmic}
\end{algorithm}

\textbf{Tier~1: Query Rewriting} ($0.5 < H < 0.65$). For mild hallucinations, the system generates three alternative query phrasings, re-retrieves context for each, deduplicates and reranks the combined chunks, and regenerates the response. This tier addresses retrieval-induced hallucination by diversifying the evidence base.

\textbf{Tier~2: Context Reranking} ($0.65 \leq H < 0.80$). For moderate hallucinations, a cross-encoder reranker (cross-encoder/ms-marco-MiniLM-L-6-v2) reorders existing context chunks by query relevance and retains only the top-3 most relevant chunks. Regeneration with higher-quality context addresses both retrieval noise and faithfulness failures.

\textbf{Tier~3: Self-Reflection} ($H \geq 0.80$). For severe hallucinations, an iterative self-correction loop prompts the LLM to identify unsupported claims in its own response, remove or correct them against the context, and regenerate. The loop runs for a maximum of two iterations or until the re-scored hallucination score drops below the threshold.

A critical design decision is the \textit{keep-best} policy (line~16 of Algorithm~\ref{alg:router}): after mitigation, the system re-scores the corrected response and returns the version with the lower hallucination score, guarding against mitigation-induced regressions.

\subsection{Claim-Level Explanation Generator}

To support human-in-the-loop verification, HalluGuard provides per-sentence explanations. The response is tokenised into sentences using NLTK's \texttt{sent\_tokenize}. Each sentence is scored individually against the context using the NLI model. For sentences with hallucination scores exceeding 0.5, the LLM generates a one-sentence natural language explanation identifying what is unsupported. The output format is: \texttt{[HALLUCINATED: \textit{claim} | Reason: \textit{explanation}]}, enabling rapid human review of flagged claims.

% ============================================================
%  IV. EXPERIMENTS
% ============================================================
\section{Experiments}

\subsection{Experimental Setup}

\textbf{Dataset.} We evaluate on the QA split of HaluEval~\cite{li2023halueval} (\texttt{pminervini/HaluEval} on HuggingFace), using 500 samples containing a balanced mix of faithful and hallucinated responses with associated context passages and ground-truth labels.

\textbf{Models.} The LLM generator is Llama-3-8B-Instruct served via the Groq API. The NLI scorer uses cross-encoder/nli-deberta-v3-base. The embedding models are sentence-transformers/all-MiniLM-L6-v2 (retrieval) and all-mpnet-base-v2 (similarity scoring). The cross-encoder reranker is cross-encoder/ms-marco-MiniLM-L-6-v2. The vector store is ChromaDB.

\textbf{Hardware.} All experiments run on an Apple M1 Air with 8\,GB RAM, using CPU only. No GPU is required---all local models are inference-optimised transformer models that execute on CPU.

\textbf{Metrics.} We report: (i)~hallucination rate (\%), the fraction of responses with $H > 0.5$; (ii)~mean faithfulness, defined as $1 - \text{mean}(H)$; (iii)~AUROC of the detector against HaluEval ground-truth labels; (iv)~F1 score at threshold $\tau = 0.5$; and (v)~latency statistics in milliseconds.

\subsection{Main Results}

Table~\ref{tab:main_results} presents the main evaluation results. On 500 HaluEval QA samples, the baseline RAG system produces a hallucination rate of 58.2\% with mean faithfulness of 0.4227. HalluGuard reduces the hallucination rate to 39.0\% and improves mean faithfulness to 0.5712, representing a 33.0\% relative reduction in hallucination and a 35.1\% relative improvement in faithfulness. The ensemble detector achieves an AUROC of 0.748 and F1 of 0.722 against HaluEval ground-truth labels at threshold $\tau = 0.5$.

\begin{table}[t]
\centering
\caption{Main Evaluation Results on HaluEval QA (500 Samples)}
\label{tab:main_results}
\begin{tabular}{lcc}
\toprule
\textbf{Metric} & \textbf{Baseline RAG} & \textbf{HalluGuard} \\
\midrule
Hallucination Rate (\%) $\downarrow$ & 58.2 & \textbf{39.0} \\
Mean Faithfulness $\uparrow$ & 0.4227 & \textbf{0.5712} \\
AUROC $\uparrow$ & 0.748 & 0.647 \\
F1 @ $\tau=0.5$ $\uparrow$ & --- & 0.722 \\
Precision @ $\tau=0.5$ & --- & 0.685 \\
Recall @ $\tau=0.5$ & --- & 0.763 \\
\midrule
No Mitigation (\%) & --- & 41.2 \\
Query Rewrite (\%) & --- & 9.8 \\
Reranking (\%) & --- & 31.8 \\
Self-Reflection (\%) & --- & 17.2 \\
\bottomrule
\end{tabular}
\end{table}

Of the 500 queries, 206 (41.2\%) received no mitigation because their hallucination score fell below the threshold, 49 (9.8\%) were routed to query rewriting, 159 (31.8\%) to context reranking, and 86 (17.2\%) to self-reflection.

The observed AUROC decrease from 0.748 (pre-mitigation) to 0.647 (post-mitigation) is an expected evaluation artefact: successfully corrected responses receive low hallucination scores from the detector, which appear as false negatives when evaluated against the original pre-correction ground-truth labels. This phenomenon is inherent to any mitigation system evaluated against static ground truth and does not represent a degradation in detection capability.

\subsection{Ablation Study}

Table~\ref{tab:ablation} presents the ablation study conducted on 200 HaluEval samples to isolate the contribution of the adaptive router versus fixed mitigation strategies.

\begin{table}[t]
\centering
\caption{Ablation Study: Fixed vs.\ Adaptive Mitigation (200 Samples)}
\label{tab:ablation}
\begin{tabular}{@{}lcccc@{}}
\toprule
\textbf{Configuration} & \textbf{H-Rate}$\downarrow$ & \textbf{Faith.}$\uparrow$ & \textbf{Lat.\ (ms)} & \textbf{Imp.\ \%} \\
\midrule
Baseline RAG & 57.0\% & 0.4312 & 104 & --- \\
NLI detector only & 57.0\% & 0.4312 & 104 & 0.0\% \\
\midrule
Fixed: Query Rewrite & 40.0\% & 0.5904 & 15{,}166 & +36.9\% \\
Fixed: Self-Reflection & 27.0\% & 0.6004 & 66{,}209 & +39.2\% \\
\midrule
\textbf{Adaptive Router} & \textbf{30.5\%} & \textbf{0.6311} & \textbf{19{,}555} & \textbf{+46.4\%} \\
\bottomrule
\end{tabular}
\end{table}

The adaptive router achieves the highest faithfulness (0.6311, +46.4\% over baseline) while maintaining moderate latency (19.6\,s mean). Fixed self-reflection achieves a lower hallucination rate (27.0\% vs.\ 30.5\%) but at 3.4$\times$ higher latency (66.2\,s) and lower faithfulness (0.6004 vs.\ 0.6311). This demonstrates that applying the most expensive strategy uniformly is suboptimal: the adaptive router achieves superior faithfulness by matching strategy intensity to hallucination severity, avoiding both undercorrection of severe cases and overcorrection of mild ones. Fixed query rewriting achieves +36.9\% faithfulness improvement but leaves 40.0\% of responses hallucinated, confirming that lightweight strategies alone are insufficient for severe hallucinations.

\subsection{Cross-Domain Generalisation}

Table~\ref{tab:cross_domain} evaluates HalluGuard across three domains to assess generalisation without domain-specific tuning.

\begin{table}[t]
\centering
\caption{Cross-Domain Generalisation Performance}
\label{tab:cross_domain}
\begin{tabular}{@{}lcccc@{}}
\toprule
\textbf{Domain} & \textbf{N} & \textbf{H-Rate}$\downarrow$ & \textbf{Faith.}$\uparrow$ & \textbf{AUROC} \\
\midrule
General (HaluEval) & 500 & 39.0\% & 0.5712 & 0.647 \\
Medical (MedQA) & 100 & 10.0\% & 0.5979 & 0.614 \\
Legal & 100 & 51.0\% & 0.4974 & 0.502 \\
\bottomrule
\end{tabular}
\end{table}

HalluGuard generalises well to the medical domain, achieving a hallucination rate of 10.0\% with faithfulness of 0.5979. The legal domain proves more challenging (51.0\% hallucination rate, AUROC 0.502), which we attribute to the higher proportion of domain-specific terminology and reasoning patterns that challenge both the NLI model and the LLM judge without domain adaptation.

\subsection{Latency Analysis}

Table~\ref{tab:latency} reports per-component latency from a 10-sample profiling run on CPU hardware.

\begin{table}[t]
\centering
\caption{Per-Component Latency Breakdown (CPU, Apple M1)}
\label{tab:latency}
\begin{tabular}{@{}lcc@{}}
\toprule
\textbf{Component} & \textbf{Mean (ms)} & \textbf{p95 (ms)} \\
\midrule
Retriever & 0.4 & 1.8 \\
Ensemble Detector (parallel) & 4{,}811 & 8{,}502 \\
Mitigation Router & 1{,}230 & 6{,}765 \\
\midrule
End-to-end (warm) & 6{,}042 & 15{,}266 \\
Cold start & \multicolumn{2}{c}{14{,}573} \\
\bottomrule
\end{tabular}
\end{table}

Parallel execution of the three detector signals reduced detection latency by 70.0\% compared to sequential execution (6{,}042\,ms vs.\ 20{,}108\,ms end-to-end). The latency distribution is heavily bimodal: queries that require no mitigation ($H \leq 0.5$, representing 41.2\% of traffic) incur only detection overhead, while queries routed to self-reflection incur substantially higher latency. This represents a configurable accuracy--efficiency trade-off rather than a fixed cost: deployments prioritising latency can cap the maximum mitigation tier, while deployments prioritising accuracy can enable all tiers with asynchronous processing for the self-reflection tier.

% ============================================================
%  V. DISCUSSION
% ============================================================
\section{Discussion}

\textbf{Why confidence-adaptive routing outperforms fixed strategies.} The ablation results (Table~\ref{tab:ablation}) reveal a fundamental insight: different hallucination severities require proportionate correction. Mild hallucinations---where the generator partially follows the context but injects minor unsupported details---are efficiently resolved by query rewriting, which diversifies the evidence base. Applying self-reflection to these cases introduces unnecessary latency (66\,s vs.\ 15\,s) without proportionate benefit. Conversely, severe hallucinations where the generator entirely ignores retrieved context cannot be resolved by merely reranking the same passages. The adaptive router's +46.4\% faithfulness improvement over baseline exceeds both fixed rewriting (+36.9\%) and fixed self-reflection (+39.2\%) by correctly matching intervention intensity to problem severity.

\textbf{Limitations.} Three limitations merit discussion. First, the self-reflection tier incurs substantial latency (mean 66\,s on CPU), making it unsuitable for sub-second applications; this is addressable through asynchronous processing where the initial response is served immediately while correction proceeds in the background. Second, cross-domain evaluation on the legal domain used a small evaluation set ($n=100$) and the AUROC of 0.502 suggests that domain-specific adaptation may be necessary for specialised domains. Third, the LLM-judge signal introduces a dependency on the Groq API; a fully offline deployment would require a locally hosted judge model.

\textbf{Comparison with Self-RAG and CRAG.} Both Self-RAG~\cite{asai2024selfrag} and CRAG~\cite{yan2024crag} apply one fixed correction strategy to all detected hallucinations. Self-RAG additionally requires supervised fine-tuning with reflection token annotations, while CRAG depends on web search for corrective retrieval. HalluGuard requires no fine-tuning, operates with local retrieval only, and selects correction strategies proportionately---a design that our ablation study demonstrates yields superior faithfulness.

% ============================================================
%  VI. CONCLUSION
% ============================================================
\section{Conclusion}

We presented HalluGuard, a training-free framework for real-time hallucination detection and mitigation in RAG systems. HalluGuard introduces four contributions: (1)~an ensemble detector combining NLI, embedding, and LLM-judge signals with parallel execution; (2)~a confidence-gated mitigation router that selects correction strategies proportionate to hallucination severity; (3)~claim-level explanations for human-in-the-loop verification; and (4)~multi-turn entity consistency tracking.

On 500 HaluEval QA samples, HalluGuard reduces the hallucination rate from 58.2\% to 39.0\% and improves mean faithfulness by 35.1\%. The adaptive router achieves +46.4\% faithfulness improvement over baseline, surpassing all fixed-strategy alternatives. The system requires no fine-tuning and operates as a plug-and-play module for any RAG pipeline.

Future work will explore asynchronous mitigation to reduce self-reflection latency, domain-specific adaptation for legal and medical corpora, and extension to multilingual RAG systems.

% ============================================================
%  REFERENCES
% ============================================================
\bibliographystyle{IEEEtran}
\bibliography{halluguard_paper}

\end{document}
